\textbf{Phần 1. Trắc nghiệm lựa chọn}
\setcounterpageref{ex}{0}

\begin{ex} %[3H13N66]
Một ngày có bao nhiêu giờ?
\choice
{$12$ giờ}
{$48$ giờ}
{$60$ giờ}
{\True $24$ giờ}
\loigiai{
Một ngày có 24 giờ.\\
\textbf{Đáp án của câu hỏi là D}
}
\end{ex}

\begin{ex} %[3H13N66]
Một giờ có bao nhiêu phút?
\choice
{$30$ phút}
{\True $60$ phút}
{$100$ phút}
{$12$ phút}
\loigiai{
Một giờ có 60 phút.\\
\textbf{Đáp án của câu hỏi là B}
}
\end{ex}

\begin{ex} %[3H13H66]
Đồng hồ chỉ kim ngắn qua số 8 một chút, kim dài chỉ số 3. Lúc đó là
\choice
{\True $8$ giờ $15$ phút}
{$8$ giờ $3$ phút}
{$8$ giờ $30$ phút}
{$3$ giờ $40$ phút}
\loigiai{
Kim dài chỉ số 3 là 15 phút $\Rightarrow 8$ giờ 15 phút.\\
\textbf{Đáp án của câu hỏi là A}
}
\end{ex}

\begin{ex} %[3H13H66]
$8$ giờ $45$ phút còn được gọi là
\choice
{$8$ giờ kém $15$ phút}
{$9$ giờ $45$ phút}
{\True $9$ giờ kém $15$ phút}
{$8$ giờ rưỡi}
\loigiai{
$8$ giờ $45$ phút $= 9$ giờ kém 15 phút.\\
\textbf{Đáp án của câu hỏi là C}
}
\end{ex}

\begin{ex} %[3H13H66]
Một năm thường có bao nhiêu tháng?
\choice
{$10$ tháng}
{$14$ tháng}
{$6$ tháng}
{\True $12$ tháng}
\loigiai{
Một năm có 12 tháng.\\
\textbf{Đáp án của câu hỏi là D}
}
\end{ex}

\begin{ex} %[3H13H66]
Các tháng có 31 ngày trong năm gồm
\choice
{Tháng 1, 2, 3, 4, 5, 6, 7}
{\True Tháng 1, 3, 5, 7, 8, 10, 12}
{Tháng 4, 6, 9, 11}
{Tháng 2, 4, 6, 8, 10, 12}
\loigiai{
Các tháng có 31 ngày: 1, 3, 5, 7, 8, 10, 12.\\
\textbf{Đáp án của câu hỏi là B}
}
\end{ex}

\begin{ex} %[3H13V66]
Tháng 4 có bao nhiêu ngày?
\choice
{$31$ ngày}
{$28$ ngày}
{\True $30$ ngày}
{$29$ ngày}
\loigiai{
Tháng 4 có 30 ngày.\\
\textbf{Đáp án của câu hỏi là C}
}
\end{ex}

\begin{ex} %[3H13V66]
Tháng 2 năm nhuận có bao nhiêu ngày?
\choice
{\True $29$ ngày}
{$28$ ngày}
{$30$ ngày}
{$31$ ngày}
\loigiai{
Năm nhuận tháng 2 có 29 ngày (năm thường có 28 ngày).\\
\textbf{Đáp án của câu hỏi là A}
}
\end{ex}

\begin{ex} %[3H13V66]
Nếu ngày 1 tháng 5 là thứ Năm thì ngày 8 tháng 5 là thứ mấy?
\choice
{$Thứ Tư$}
{$Thứ Sáu$}
{\True $Thứ Năm$}
{$Thứ Bảy$}
\loigiai{
Cách 7 ngày (1 tuần) nên vẫn là Thứ Năm ($1 + 7 = 8$).\\
\textbf{Đáp án của câu hỏi là C}
}
\end{ex}

\begin{ex} %[3H13V66]
Lan bắt đầu học bài lúc 8 giờ tối và học xong lúc 9 giờ kém 15 phút tối. Lan đã học bài trong bao nhiêu phút?
\choice
{$40$ phút}
{\True $45$ phút}
{$50$ phút}
{$60$ phút}
\loigiai{
$9$ giờ kém 15 phút $= 8$ giờ 45 phút $\Rightarrow$ Học trong 45 phút.\\
\textbf{Đáp án của câu hỏi là B}
}
\end{ex}

\begin{ex} %[3H13V66]
Các tháng nào sau đây có đúng 30 ngày?
\choice
{Tháng 1, 3, 5, 7}
{Tháng 7, 8, 10, 12}
{Tháng 2, 4, 6, 8}
{\True Tháng 4, 6, 9, 11}
\loigiai{
Các tháng có 30 ngày: 4, 6, 9, 11.\\
\textbf{Đáp án của câu hỏi là D}
}
\end{ex}

\begin{ex} %[3H13C66]
Ngày 28 tháng 4 là thứ Ba. Hỏi ngày 1 tháng 5 cùng năm đó là thứ mấy?
\choice
{\True $Thứ Sáu$}
{$Thứ Năm$}
{$Thứ Bảy$}
{$Chủ Nhật$}
\loigiai{
Tháng 4 có 30 ngày. Ngày 28 là Thứ Ba $\rightarrow$ Ngày 29 (Thứ Tư) $\rightarrow$ Ngày 30 (Thứ Năm) $\rightarrow$ Ngày 1 tháng 5 là Thứ Sáu.\\
\textbf{Đáp án của câu hỏi là A}
}
\end{ex}

\textbf{Phần 2. Trắc nghiệm đúng sai}
\setcounterpageref{ex}{0}

\begin{ex} %[3H13H66]
Xét tính đúng sai về thời gian:
\choiceTF
{\True Một năm có 12 tháng.}
{\True Tháng 7 và tháng 8 đều có 31 ngày liên tiếp.}
{Tháng 6 có 31 ngày.}
{\True Tháng 6 có đúng 30 ngày.}
\loigiai{
\begin{itemize}
    \item Một năm có 12 tháng: Đúng.
    \item Tháng 7 và tháng 8 đều có 31 ngày: Đúng.
    \item Tháng 6 có 30 ngày (không phải 31). Ý 3 Sai, Ý 4 Đúng.
\end{itemize}
\begin{itemize}
    \item a) \textbf{Đúng}.
    \item b) \textbf{Đúng}.
    \item c) \textbf{Sai}.
    \item d) \textbf{Đúng}.
\end{itemize}
}
\end{ex}

\begin{ex} %[3H13V66]
Đồng hồ chỉ 3 giờ 15 phút chiều (tức 15 giờ 15 phút). Xét tính đúng sai:
\choiceTF
{\True Kim ngắn chỉ qua số 3 một chút.}
{\True Kim dài chỉ đúng số 3 (15 phút).}
{Lúc đó là 15 giờ 3 phút.}
{\True Lúc đó là 15 giờ 15 phút.}
\loigiai{
\begin{itemize}
    \item Kim ngắn qua số 3: Đúng.
    \item Kim dài chỉ số 3 (15 phút): Đúng.
    \item 15 giờ 15 phút (không phải 3 phút). Ý 3 Sai, Ý 4 Đúng.
\end{itemize}
\begin{itemize}
    \item a) \textbf{Đúng}.
    \item b) \textbf{Đúng}.
    \item c) \textbf{Sai}.
    \item d) \textbf{Đúng}.
\end{itemize}
}
\end{ex}

\textbf{Phần 3. Trả lời ngắn}
\setcounterpageref{ex}{0}

\begin{ex} %[3H13H66]
Tháng 12 có bao nhiêu ngày?
\shortans{31}
\loigiai{
\textbf{Đáp số:} $31$\\
31 ngày.
}
\end{ex}

\begin{ex} %[3H13H66]
Tháng 9 có bao nhiêu ngày?
\shortans{30}
\loigiai{
\textbf{Đáp số:} $30$\\
30 ngày.
}
\end{ex}

\begin{ex} %[3H13V66]
Một tiết học bắt đầu lúc 8 giờ và kết thúc lúc 8 giờ 35 phút. Tiết học đó kéo dài bao nhiêu phút?
\shortans{35}
\loigiai{
\textbf{Đáp số:} $35$\\
35 phút.
}
\end{ex}

\begin{ex} %[3H13V66]
Có bao nhiêu tháng trong năm có đúng 31 ngày?
\shortans{7}
\loigiai{
\textbf{Đáp số:} $7$\\
Có 7 tháng (1, 3, 5, 7, 8, 10, 12).
}
\end{ex}

\begin{ex} %[3H13V66]
Nếu ngày 2 tháng 9 là thứ Ba thì ngày 9 tháng 9 là thứ mấy (chỉ ghi số thứ: ví dụ thứ Ba ghi 3)?
\shortans{3}
\loigiai{
\textbf{Đáp số:} $3$\\
$2 + 7 = 9$ nên vẫn là thứ Ba (ghi 3).
}
\end{ex}

\begin{ex} %[3H13C66]
Năm không nhuận có tất cả bao nhiêu ngày?
\shortans{365}
\loigiai{
\textbf{Đáp số:} $365$\\
365 ngày.
}
\end{ex}
